\providecommand{\pdfxopt}{a-1b,mathxmp}

%\documentclass[10pt]{article}
%\usepackage{beamerarticle}
\documentclass[ignorenonframetext,aspectratio=169]{beamer}

\input{/home/asreeve/current/tex_style/preamble_102}
\usepackage{cancel}

\setbeamertemplate{enumerate item}[square]
\setbeamertemplate{enumerate subitem}[circle]
\setbeamertemplate{enumerate subsubitem}[square]
\setbeamercolor*{enumerate item}{fg=red,bg=blue}
\setbeamercolor*{enumerate subitem}{fg=blue,bg=red}

\begin{document}

\begin{frame}{Vernal Pools}
  \begin{minipage}{.68\textwidth}
  Vernal pools, a wetland system found in Maine and many other locations, form in wet periods and then slowly dry out over time.
  \begin{enumerate}
  \item \small These pools occur in isolated depressions with no channelized input or output.
  \item \small Vernal pools are important habitat for a variety of organisms, especially amphibians that can thrive in these fish-less pools.
  \item \small They export nutrients and food (delicious frogs) to the upland.
  \end{enumerate}

  In this lab, you will evaluate the hydrology of vernal pools.
\end{minipage}
\begin{minipage}{.3\textwidth}
  \includegraphics[width=\textwidth]{../../figs/20201007_191023_crop.jpg}
  \includegraphics[width=\textwidth]{../../figs/20160507_152637.jpg}
  \tiny Spotted salamander (top) and egg masses in vernal pool.
\end{minipage}

\end{frame}

\begin{frame}{The Water Balance}
  \begin{minipage}{.30\textwidth}
    \includegraphics[width=\textwidth]{../../figs/DSC01850.JPG}
    \includegraphics[width=\textwidth]{../../figs/dscn0248.jpg}
    \tiny Green frog in vernal pool(top) and future earth scientist pointing to vernal pool. 
  \end{minipage}
  \begin{minipage}{.68\textwidth}
    \begin{itemize}
    \item \small A water budget tracks the inputs and outputs of water and compares this to water volumes measured in a reservoir:
      \begin{equation}
        \Delta S = Vol_{in}-Vol_{out}
      \end{equation}
      \tiny measured over some time interval \normalsize
      
      or
      \begin{equation}
        \Delta S = \Delta t \left( Flux_{in}-Flux_{out} \right)
      \end{equation}
    \item \small Comparing each side of the equation separately allows errors or problems in the methods can be identified.
    \item There are many types of fluxes that can impact a water body.
    \end{itemize}
  \end{minipage}    
\end{frame}

\begin{frame}{Calculating a balance with a spreadsheet.}
  \begin{minipage}{.3\textwidth}
    \includegraphics[width=\textwidth]{../../figs/dscn0526.jpg}
    \includegraphics[width=\textwidth]{../../figs/dscn0572.jpg}
    \tiny A wood frog at the edge of a wetland (top) and green frog in vernal pool (bottom)
  \end{minipage}
  \begin{minipage}{.68\textwidth}
    A spreadsheet can be used to do the calculations needed for a water budget. Using data collected from vernal pools in Osborn and Bangor Maine, construct a water budget for these very different vernal pools and use it to consider the processes that are most important in a vernal pool.

    The next slides will review the use of a spreadsheet to do these calculation. The intent here is to teach you how to use a spreadsheet so you can apply it to other problems. The \textbf{example data and descriptions will not necessarily match the problem presented in the lab}. Do not simply mimic what is presented here, but instead apply the concepts to the problem. I will use google sheets in the example; you are free to use another spreadsheet program to complete the lab assignment.
  \end{minipage}
\end{frame}
  
\end{document}
